\chapter{John F. Nash Jr. (2015)}\index{Nash Jr., John F.}

\section{Biographical Sketch}

John Forbes Nash Jr. was born on 13 June 1928 in Bluefield, West Virginia, USA. He studied at Carnegie Institute of Technology (now Carnegie Mellon University) and Princeton University, where he completed his PhD in 1950 under the supervision of Albert W. Tucker. His 27-page doctoral dissertation introduced the concept of the Nash equilibrium\index{Nash equilibrium}, which revolutionized economics. Nash held positions at the Massachusetts Institute of Technology and Princeton University. He received the Nobel Prize in Economics in 1994 and the Abel Prize (jointly with Louis Nirenberg) in 2015. Nash's life and struggle with schizophrenia were the subject of the book and film ``A Beautiful Mind.''

\section{High-Level View for a General Audience}

\subsection{What Did Nash Do?}

Think about the shape of a soap film stretched across a wire frame. The film takes the shape that minimizes its surface area. This is a problem of geometry: given a boundary, find the surface of least area spanning it.

Now imagine that you are told the curvature of a surface at every point, but not its overall shape. Can you reconstruct the shape from this curvature information? This is the kind of problem that John Nash solved.

Nash proved that any Riemannian manifold can be isometrically embedded\index{Isometric embedding} into Euclidean space with sufficiently many dimensions. In simple terms: any curved space can be placed inside a flat space of higher dimension in a way that preserves all distances. This is the Nash embedding theorem, one of the most surprising results in geometry.

\subsection{Why Does It Matter?}

\begin{itemize}
\item \textbf{Geometry and Physics}: The Nash embedding theorem tells us that curved spaces (like spacetime in general relativity) can be studied as subsets of flat spaces of higher dimension. This is a key insight for both mathematics and physics.

\item \textbf{Materials Science}: The mathematics of surfaces and their curvatures, which Nash studied, is used to understand the behavior of thin films, crystals, and other materials.

\item \textbf{Animation and Graphics}: The reconstruction of surfaces from geometric data, which uses the mathematics of Nash, is used in computer graphics and animation.

\item \textbf{Economics}: While separately famous for Nash equilibrium\index{Nash equilibrium} in game theory\index{Game theory}, Nash's mathematical work on embedding theorems has deep connections to the geometry of economic models and their representations.
\end{itemize}

\section{The Origin of the Problem}

\subsection{The Embedding Problem}

The embedding problem asks: given an abstract Riemannian manifold\index{Manifolds} $(M^n, g)$, can it be realized as a submanifold of Euclidean space $\mathbb{R}^N$ in such a way that the induced metric equals $g$?

The isometric embedding problem was posed by Riemann himself in his 1854 habilitation lecture. By the 1950s, the following was known:
\begin{itemize}
\item Local isometric embeddings exist for analytic metrics (Cauchy--Kovalevskaya theorem).
\item The Whitney embedding theorem provides topological embeddings into $\mathbb{R}^{2n}$.
\item But isometric embeddings---those that preserve distances---seemed much harder.
\end{itemize}

\begin{knowledgebridge}{What Is an Isometric Embedding?}
An isometric embedding places a curved space inside a flat space without stretching or tearing---preserving all internal distances. Think of crumpling a piece of paper: the paper's intrinsic geometry (distances measured along the surface) changes. An isometric embedding is like gently draping fabric over a curved frame so that the fabric's own measurements are preserved. Nash proved that any curved space, no matter how abstract, can be placed inside a flat Euclidean space of sufficiently high dimension with all distances intact.
\end{knowledgebridge}

\section{Deep Insight into the Problem}

\subsection{Nash's \texorpdfstring{$C^1$}{C\^{}1} Embedding Theorem}

In 1954, Nash stunned the mathematical world by proving that any Riemannian manifold can be isometrically embedded into Euclidean space. His first proof (the $C^1$ embedding theorem) used a remarkably simple idea: a series of small perturbations that gradually reduce the error in the embedding.

The $C^1$ Nash embedding theorem: every compact Riemannian $n$-manifold admits a $C^1$ isometric embedding into $\mathbb{R}^{n+1}$.

The construction works as follows:
\begin{enumerate}
\item Start with a short embedding (one that shrinks all distances).
\item Decompose the difference between the metric $g$ and the pullback of the Euclidean metric into a sum of small ``elementary'' metrics.
\item Each elementary metric corresponds to a perturbation that adds a wrinkle to the embedding in a new direction.
\item With infinitely many wrinkles in $n+1$ directions, the metric error is corrected without changing the $C^1$ topology.
\end{enumerate}

\begin{bigidea}
Nash's revolutionary idea for $C^1$ embeddings is that adding tiny wrinkles in a new direction can correct metric errors without creating self-intersections. Start with an embedding that is too short (all distances are compressed). Decompose the metric error into a sum of small pieces, each of which can be fixed by adding a zigzag in a new dimension. By using infinitely many dimensions (directions), Nash can make the error vanish while keeping the map only $C^1$---continuous with a derivative, but not twice differentiable. The wrinkles do the work.
\end{bigidea}

\subsection{The Smooth Nash Embedding Theorem}

The smooth version requires more dimensions. Nash proved:
\[
(M^n, g) \text{ embeds isometrically into } \mathbb{R}^{n(3n+11)/2}.
\]

The dimension bound was later improved by Gromov and Gunther to $\frac{n(n+1)}{2} + 2n + 3$ and then to $\max\{ n(n+1)/2 + n, n(n+1)/2 + 5 \}$.

\subsection{The Nash--Moser Implicit Function Theorem}

Nash's proof of the smooth embedding theorem required a new kind of implicit function theorem, now called the Nash--Moser theorem. Unlike the classical inverse function theorem, which works in Banach spaces, the Nash--Moser theorem works in Fr\'echet spaces and requires smoothing operations at each step of the iteration.

The technical innovation: the iteration scheme includes a smoothing operator $S_\theta$ that regularizes the solution at each step. The ``tame'' estimate ensures that the smoothing does not amplify errors. The theorem was refined by Moser (hence the name Nash--Moser) and later by Gromov and Hamilton.

\subsection{The Nash Twist}

The key geometric construction in Nash's proof is the ``Nash twist''---a periodic deformation of the embedding that adds a short-wavelength correction to the metric without creating self-intersections. The twist is controlled by a rapidly oscillating function, and the product structure ensures that the different corrections do not interfere.

\begin{analogy}{The Tailor's Mannequin}
Nash's twist works like a tailor adjusting a suit: to make a sleeve longer without tearing the fabric, the tailor adds a series of tiny pleats that expand the effective length. Each Nash twist adds a fast ripple to the embedding in a new direction, stretching the metric slightly in that direction without affecting the other directions. The ripples are so rapid that they average out smoothly, and because each twist operates in its own coordinate direction, none of the corrections interfere with each other.
\end{analogy}

\section{Biggest Contributions and New Tools}

\subsection{The Nash Embedding Theorem}

Nash's embedding theorem comes in two versions:
\[
\text{Every smooth Riemannian manifold } (M^n, g) \text{ embeds isometrically into } \mathbb{R}^{n(3n+11)/2}.
\]

The proof introduced the ``Nash twist'' and the iteration method that came to be called Nash--Moser iteration.

\subsection{The Nash--Moser Theorem}

This theorem provides an inverse function theorem in Fr\'echet spaces under the conditions of a ``tame'' estimate and a smooth approximation scheme. The theorem is essential for Nash's smooth embedding theorem and has applications in KAM theory and the geometry of diffeomorphism groups.

\section{Impact of the Discovery}

\subsection{Impact on Mathematics}

\begin{itemize}
\item \textbf{Differential Geometry}: The Nash embedding theorem is a fundamental result that justifies the study of Riemannian manifolds as subsets of Euclidean space.

\item \textbf{Functional Analysis}: The Nash--Moser theorem is a powerful tool in the analysis of nonlinear PDEs\index{PDEs} in geometry.

\end{itemize}

\begin{whyitmatters}
Nash's embedding theorem reshaped how mathematicians think about geometry: every abstract Riemannian manifold can be realized as a concrete submanifold of Euclidean space. This means that the entire apparatus of Euclidean geometry---angles, distances, curvatures computed from an ambient coordinate system---is universally available.
\end{whyitmatters}

\subsection{Impact on Technology}

\begin{itemize}
\item \textbf{Materials Science}: The geometry of thin films and interfaces uses the theory of surfaces and curvature that Nash advanced.
\end{itemize}

\subsection{Impact on Economics}

Nash's equilibrium concept, now called ``Nash equilibrium,'' transformed economics and social science. It provides the mathematical foundation for analyzing strategic interactions in markets, politics, and social behavior.

\section{Current Developments}

\subsection{Isometric Embedding in Low Codimensions}

The minimal dimension needed for isometric embeddings remains an active area of research. Gunther improved Nash's bound to $N = \max\{n(n+1)/2 + n, n(n+1)/2 + 5\}$. The optimal dimension is conjectured to be $n(n+1)/2$ for compact manifolds and $n(n+1)/2 + n$ for non-compact ones.

\subsection{Fluid Dynamics and the Nash--Moser Theorem}

The Nash--Moser theorem has been applied to the study of the Euler equations. In particular, De Lellis and Sz\'ekelyhidi used Nash-type iteration to construct non-unique weak solutions of the Euler equations, establishing the failure of Onsager's conjecture for low regularity.

\subsection{Game Theory and Machine Learning}

Nash equilibrium concepts are used in multi-agent reinforcement learning and the design of machine learning algorithms. The theory of generative adversarial networks (GANs) is based on a two-player game, and the training dynamics are analyzed using game theory.

\subsection{The Isometric Embedding of the Sphere}

The isometric embedding of the sphere $(S^2, g_0)$ into $\mathbb{R}^3$ is a classical problem. While $S^2$ cannot be isometrically embedded into $\mathbb{R}^2$, the $C^1$ embedding theorem says it can be embedded into $\mathbb{R}^3$ with a $C^1$ but not $C^2$ map (this is the subject of the Nash--Kuiper theorem, which constructs $C^1$ isometric embeddings of surfaces into $\mathbb{R}^3$ that are arbitrarily close to any short embedding).

\section{Conclusion}

John Nash transformed our understanding of differential geometry. His embedding theorem is one of the most surprising results in geometry, showing that any curved space can be realized as a subset of flat space. Together with Louis Nirenberg, he received the Abel Prize ``for striking and seminal contributions to the theory of nonlinear partial differential equations\index{PDEs} and its applications to geometric analysis.''

\section{Behind the Breakthrough: The Human Story}
Nash's embedding theorem was born from a bet. A colleague claimed Nash couldn't prove that any Riemannian manifold could be embedded in Euclidean space. Nash went to work and proved it, proving the colleague wrong.

\section{Pedagogical Metaphors and Lineage}
\subsection{Signature Metaphor}
``Folding and weaving any curved space smoothly into Euclidean space without stretching it, like wrapping a gift without creasing the paper.''

\subsection{Mathematical Lineage}
Nash advised by Albert W. Tucker.

\section{Applied Science and Computational Algorithms}
\subsection{Key Takeaways for Applied Science}
Applied in computer graphics (isometric texture mapping, mesh deformation, and skinning) and in aerospace engineering to solve fluid flow and heat dissipation boundaries.

\subsection{Algorithms and Numerical Implementations}
Least Squares Conformal Maps (LSCM) and geometric flattening algorithms in CAD software solve boundary value problems using Nash's embedding coordinates.

\section{The Research Frontier}
Proving global regularity or blow-up of solutions to the 3D Navier-Stokes equations, which is a major open problem in mathematical physics.

\section{Guided Exercises and Challenges}
\begin{enumerate}[label=\textbf{\arabic*.}]
  \item State the Nash embedding theorem and explain the difference between $C^1$ isometric embeddings and smooth $C^k$ embeddings.
  \item Explain the Nash-Moser inverse function theorem and why it is needed when solving non-linear PDEs where derivative loss occurs.
  \item State the De Giorgi-Nash-Moser theorem on the H\"older continuity of weak solutions to elliptic equations.
\end{enumerate}

\section{Curriculum Map and Further Reading}
For students wishing to delve deeper into the mathematical machinery underlying these breakthroughs, the following learning path is recommended:
\begin{itemize}
  \item John Nash, \emph{The Imbedding Problem for Riemannian Manifolds}, Annals of Mathematics 63 (1956).
  \item L. Nirenberg, \emph{Topics in Nonlinear Functional Analysis}, AMS.
\end{itemize}

\section*{Bibliography}
\begin{enumerate}
\item J. F. Nash Jr., ``$C^1$ isometric imbeddings,'' Annals of Mathematics 60 (1954), 383--396.
\item J. F. Nash Jr., ``The imbedding problem for Riemannian manifolds,'' Annals of Mathematics 63 (1956), 20--63.
\item J. F. Nash Jr., ``Non-cooperative games,'' Annals of Mathematics 54 (1951), 286--295.
\item J. F. Nash Jr., ``Equilibrium points in $n$-person games,'' Proceedings of the National Academy of Sciences 36 (1950), 48--49.
\item R. S. Hamilton, ``The inverse function theorem of Nash and Moser,'' Bulletin of the American Mathematical Society 7 (1982), 65--222.
\end{enumerate}
